% Generated by build_jcaa_bibliography.py from references.bib -- do not edit by hand.
\begingroup\raggedright
\begin{thebibliography}{29}

\bibitem[Allouche, Tsoar \& Kadmon(2006)]{allouche2006}
Allouche, O, Tsoar, A and Kadmon, R. 2006. Assessing the accuracy of species distribution models: prevalence, kappa and the true skill statistic (TSS). \textit{Journal of Applied Ecology}, 43(6): 1223–1232. DOI: 10.1111/j.1365-2664.2006.01214.x

\bibitem[Ara{\'u}jo \& Guisan(2006)]{araujo2006}
Ara{\'u}jo, MB and Guisan, A. 2006. Five (or so) challenges for species distribution modelling. \textit{Journal of Biogeography}, 33(10): 1677–1688. DOI: 10.1111/j.1365-2699.2006.01584.x

\bibitem[Banks et~al.(2006)]{banks2006}
Banks, WE, d'Errico, F, Dibble, HL, Krishtalka, L, West, D, Olszewski, DI, Peterson, AT, Anderson, DG, Gillam, JC, Montet-White, A, Crucifix, M, Marean, CW, S{\'a}nchez-Go{\~n}i, M-F, Wohlfarth, B and Vanhaeren, M. 2006. Eco-cultural niche modeling: new tools for reconstructing the geography and ecology of past human populations. \textit{PaleoAnthropology}, 68–83.

\bibitem[Barbet-Massin et~al.(2012)]{barbet2012}
Barbet-Massin, M, Jiguet, F, Albert, CH and Thuiller, W. 2012. Selecting pseudo-absences for species distribution models: how, where and how many? \textit{Methods in Ecology and Evolution}, 3(2): 327–338. DOI: 10.1111/j.2041-210X.2011.00172.x

\bibitem[Barve et~al.(2011)]{barve2011}
Barve, N, Barve, V, Jim{\'e}nez-Valverde, A, Lira-Noriega, A, Maher, SP, Peterson, AT, Sober{\'o}n, J and Villalobos, F. 2011. The crucial role of the accessible area in ecological niche modeling and species distribution modeling. \textit{Ecological Modelling}, 222(11): 1810–1819. DOI: 10.1016/j.ecolmodel.2011.02.011

\bibitem[Breiman(2001)]{breiman2001}
Breiman, L. 2001. Random forests. \textit{Machine Learning}, 45(1): 5–32. DOI: 10.1023/A:1010933404324

\bibitem[Castiello \& Tonini(2021)]{castiello2021}
Castiello, ME and Tonini, M. 2021. An explorative application of Random Forest algorithm for archaeological predictive modeling. A Swiss case study. \textit{Journal of Computer Applications in Archaeology}, 4(1): 110–125. DOI: 10.5334/jcaa.71

\bibitem[Chen \& Guestrin(2016)]{chen2016}
Chen, T and Guestrin, C. 2016. XGBoost: A scalable tree boosting system. In: \textit{Proceedings of the 22nd ACM SIGKDD International Conference on Knowledge Discovery and Data Mining}. ACM, pp. 785–794. DOI: 10.1145/2939672.2939785

\bibitem[Comer et~al.(2023)]{comer2023}
Comer, JA, Comer, DC, Dumitru, IA, Priebe, CE and Patsolic, JL. 2023. Sampling methods for archaeological predictive modeling: Spatial autocorrelation and model performance. \textit{Journal of Archaeological Science: Reports}, 48: Article 103824. DOI: 10.1016/j.jasrep.2022.103824

\bibitem[Elith et~al.(2006)]{elith2006}
Elith, J, Graham, CH, Anderson, RP, Dud{\'{\i}}k, M, Ferrier, S, Guisan, A, Hijmans, RJ, Huettmann, F, Leathwick, JR, Lehmann, A, Li, J, Lohmann, LG, Loiselle, BA, Manion, G, Moritz, C, Nakamura, M, Nakazawa, Y, Overton, JMcC, Peterson, AT, Phillips, SJ, Richardson, K, Scachetti-Pereira, R, Schapire, RE, Sober{\'o}n, J, Williams, S, Wisz, MS and Zimmermann, NE. 2006. Novel methods improve prediction of species' distributions from occurrence data. \textit{Ecography}, 29(2): 129–151. DOI: 10.1111/j.2006.0906-7590.04596.x

\bibitem[Franklin(2009)]{franklin2009}
Franklin, J. 2009. \textit{Mapping Species Distributions: Spatial Inference and Prediction.} Cambridge: Cambridge University Press. DOI: 10.1017/CBO9780511810602

\bibitem[Guillera-Arroita et~al.(2015)]{guillera2015}
Guillera-Arroita, G, Lahoz-Monfort, JJ, Elith, J, Gordon, A, Kujala, H, Lentini, PE, McCarthy, MA, Tingley, R and Wintle, BA. 2015. Is my species distribution model fit for purpose? Matching data and models to applications. \textit{Global Ecology and Biogeography}, 24(3): 276–292. DOI: 10.1111/geb.12268

\bibitem[Hijmans(2012)]{hijmans2012}
Hijmans, RJ. 2012. Cross-validation of species distribution models: removing spatial sorting bias and calibration with a null model. \textit{Ecology}, 93(3): 679–688. DOI: 10.1890/11-0826.1

\bibitem[Howey, Palace \& McMichael(2016)]{howey2016}
Howey, MCL, Palace, MW and McMichael, CH. 2016. Geospatial modeling approach to monument construction using Michigan from A.D. 1000--1600 as a case study. \textit{Proceedings of the National Academy of Sciences}, 113(27): 7443–7448. DOI: 10.1073/pnas.1603450113

\bibitem[Kvamme(2006)]{kvamme2006}
Kvamme, KL. 2006. There and back again: Revisiting archaeological locational modeling. In: Mehrer, MW and Wescott, KL (eds.) \textit{GIS and Archaeological Site Location Modeling}. Boca Raton: Taylor \& Francis, pp. 3–38. DOI: 10.1201/9780203563359.sec1

\bibitem[Lobo, Jim{\'e}nez-Valverde \& Hortal(2010)]{lobo2010}
Lobo, JM, Jim{\'e}nez-Valverde, A and Hortal, J. 2010. The uncertain nature of absences and their importance in species distribution modelling. \textit{Ecography}, 33(1): 103–114. DOI: 10.1111/j.1600-0587.2009.06039.x

\bibitem[Lobo, Jim{\'e}nez-Valverde \& Real(2008)]{lobo2008}
Lobo, JM, Jim{\'e}nez-Valverde, A and Real, R. 2008. AUC: a misleading measure of the performance of predictive distribution models. \textit{Global Ecology and Biogeography}, 17(2): 145–151. DOI: 10.1111/j.1466-8238.2007.00358.x

\bibitem[Noviello et~al.(2018)]{noviello2018}
Noviello, M, Cafarelli, B, Calculli, C, Sarris, A and Mairota, P. 2018. Investigating the distribution of archaeological sites: multiparametric vs probability models and potentials for remote sensing data. \textit{Applied Geography}, 95: 34–44. DOI: 10.1016/j.apgeog.2018.04.005

\bibitem[Phillips, Anderson \& Schapire(2006)]{phillips2006}
Phillips, SJ, Anderson, RP and Schapire, RE. 2006. Maximum entropy modeling of species geographic distributions. \textit{Ecological Modelling}, 190(3–4): 231–259. DOI: 10.1016/j.ecolmodel.2005.03.026

\bibitem[Phillips et~al.(2009)]{phillips2009}
Phillips, SJ, Dud{\'{\i}}k, M, Elith, J, Graham, CH, Lehmann, A, Leathwick, J and Ferrier, S. 2009. Sample selection bias and presence-only distribution models: implications for background and pseudo-absence data. \textit{Ecological Applications}, 19(1): 181–197. DOI: 10.1890/07-2153.1

\bibitem[Radosavljevic \& Anderson(2014)]{radosavljevic2014}
Radosavljevic, A and Anderson, RP. 2014. Making better Maxent models of species distributions: complexity, overfitting and evaluation. \textit{Journal of Biogeography}, 41(4): 629–643. DOI: 10.1111/jbi.12227

\bibitem[Roberts et~al.(2017)]{roberts2017}
Roberts, DR, Bahn, V, Ciuti, S, Boyce, MS, Elith, J, Guillera-Arroita, G, Hauenstein, S, Lahoz-Monfort, JJ, Schr{\"o}der, B, Thuiller, W, Warton, DI, Wintle, BA, Hartig, F and Dormann, CF. 2017. Cross-validation strategies for data with temporal, spatial, hierarchical, or phylogenetic structure. \textit{Ecography}, 40(8): 913–929. DOI: 10.1111/ecog.02881

\bibitem[Senay, Worner \& Ikeda(2013)]{senay2013}
Senay, SD, Worner, SP and Ikeda, T. 2013. Novel three-step pseudo-absence selection technique for improved species distribution modelling. \textit{PLOS ONE}, 8(8): Article e71218. DOI: 10.1371/journal.pone.0071218

\bibitem[Valavi et~al.(2019)]{valavi2019}
Valavi, R, Elith, J, Lahoz-Monfort, JJ and Guillera-Arroita, G. 2019. blockCV: An R package for generating spatially or environmentally separated folds for k-fold cross-validation of species distribution models. \textit{Methods in Ecology and Evolution}, 10(2): 225–232. DOI: 10.1111/2041-210X.13107

\bibitem[Verhagen \& Whitley(2012)]{verhagen2012}
Verhagen, P and Whitley, TG. 2012. Integrating archaeological theory and predictive modeling: a live report from the scene. \textit{Journal of Archaeological Method and Theory}, 19(1): 49–100. DOI: 10.1007/s10816-011-9102-7

\bibitem[Wang, Shi \& Oguchi(2023)]{wang2023}
Wang, Y, Shi, X and Oguchi, T. 2023. Archaeological predictive modeling using machine learning and statistical methods for Japan and China. \textit{ISPRS International Journal of Geo-Information}, 12(6): Article 238. DOI: 10.3390/ijgi12060238

\bibitem[Warren \& Seifert(2011)]{warren2011}
Warren, DL and Seifert, SN. 2011. Ecological niche modeling in Maxent: the importance of model complexity and the performance of model selection criteria. \textit{Ecological Applications}, 21(2): 335–342. DOI: 10.1890/10-1171.1

\bibitem[Wisz \& Guisan(2009)]{wisz2009}
Wisz, MS and Guisan, A. 2009. Do pseudo-absence selection strategies influence species distribution models and their predictions? An information-theoretic approach based on simulated data. \textit{BMC Ecology}, 9: Article 8. DOI: 10.1186/1472-6785-9-8

\bibitem[Yaworsky et~al.(2020)]{yaworsky2020}
Yaworsky, PM, Vernon, KB, Spangler, JD, Brewer, SC and Codding, BF. 2020. Advancing predictive modeling in archaeology: an evaluation of regression and machine learning methods on the Grand Staircase-Escalante National Monument. \textit{PLOS ONE}, 15(10): Article e0239424. DOI: 10.1371/journal.pone.0239424

\end{thebibliography}
\endgroup
